\begin{figure}[htbp]
\centering
\setlength{\tabcolsep}{1pt}
\renewcommand{\arraystretch}{1.05}

\newcommand{\poetryimg}[1]{%
\includegraphics[
  width=0.235\linewidth,
  height=0.20\textheight,
  keepaspectratio
]{#1}
}

\begin{tabular}{cccc}

\multicolumn{1}{c}{\textbf{\footnotesize Pop Art}} &
\multicolumn{1}{c}{\textbf{\footnotesize 3D Animated Film}} &
\multicolumn{1}{c}{\textbf{\footnotesize Comic}} &
\multicolumn{1}{c}{\textbf{\footnotesize Chinese Classical Painting}} \\[6pt]

\poetryimg{figures/sec07_application/chinese_poetry/0118_8500_jpg_prompt1_img3.jpg} &
\poetryimg{figures/sec07_application/chinese_poetry/0118_8500_jpg_prompt2_img3.jpg} &
\poetryimg{figures/sec07_application/chinese_poetry/0118_8500_jpg_prompt3_img4.jpg} &
\poetryimg{figures/sec07_application/chinese_poetry/0118_8500_jpg_prompt4_img1.jpg} \\
\multicolumn{4}{c}{\footnotesize (a) \zh{月到中秋例属苏}} \\[10pt]

\poetryimg{figures/sec07_application/chinese_poetry/0118_8500_jpg_prompt21_img1.jpg} &
\poetryimg{figures/sec07_application/chinese_poetry/0118_8500_jpg_prompt22_img2.jpg} &
\poetryimg{figures/sec07_application/chinese_poetry/0118_8500_jpg_prompt23_img3.jpg} &
\poetryimg{figures/sec07_application/chinese_poetry/0118_8500_jpg_prompt24_img2.jpg} \\
\multicolumn{4}{c}{\footnotesize (b) \zh{曾经沧海难为水，从此桃源便是家}} \\[10pt]

\poetryimg{figures/sec07_application/chinese_poetry/c_0.jpg} &
\poetryimg{figures/sec07_application/chinese_poetry/c_1.jpg} &
\poetryimg{figures/sec07_application/chinese_poetry/c_2.jpg} &
\poetryimg{figures/sec07_application/chinese_poetry/c_3.jpg} \\
\multicolumn{4}{c}{\footnotesize (c) \zh{几千年兴与亡，人间正道是沧桑}} \\[10pt]


\poetryimg{figures/sec07_application/chinese_poetry/0118_8500_jpg_prompt37_img3.jpg} &
\poetryimg{figures/sec07_application/chinese_poetry/0118_8500_jpg_prompt38_img3.jpg} &
\poetryimg{figures/sec07_application/chinese_poetry/0118_8500_jpg_prompt39_img2.jpg} &
\poetryimg{figures/sec07_application/chinese_poetry/0118_8500_jpg_prompt40_img1.jpg} \\
\multicolumn{4}{c}{\footnotesize (d) \zh{梦凝白阑干，化为飞雾}} \\[6pt]

\end{tabular}

\caption{Visualization of images generated directly from classical Chinese poetry by our distilled 4-step text-to-image model under diverse artistic styles. Each row corresponds to a representative poetic verse, while each column presents the generated result in a
different style, including Pop Art, 3D Animated Film, Comic, and Chinese Classical Painting. The results show that the model is able to understand and render the semantic content, visual imagery, and aesthetic mood of classical Chinese poetry with high fidelity and
strong stylistic controllability, without requiring external English translation.}
\label{fig:chinese-poetry}
\end{figure}

\section{Application: Classical Chinese Poetry Image Generation}

To demonstrate the practical applicability of JuZhou 1.0 on a real-world culturally grounded task, we developed a mobile application named Mojie (\zh{墨界}). The application is built upon our poetry-oriented generative model, which is fine-tuned on a large-scale classical Chinese poetry dataset constructed through our own curation pipeline. After fine-tuning, the model is deployed as an Android edge application, enabling users to access poetry-to-image generation capabilities directly on mobile devices. The Android application is publicly available for download.\footnote{\url{https://www.pgyer.com/mojiemobilellm-android}}

\subsection{Application Overview}

Mojie is a practical mobile system that transforms classical Chinese poetry into high-quality images. By deploying the fine-tuned model directly on Android devices rather than relying on cloud inference, it demonstrates user-facing on-device generation in real-world applications.

The underlying model is fine-tuned on our synthetic poem-image corpus to better capture poetic semantics, scene composition, and traditional Chinese artistic imagery. Through this domain adaptation, Mojie generates visually coherent and semantically aligned images from poetic inputs. As an end-to-end prototype, it integrates data curation, model fine-tuning, edge deployment, and user interaction. Fig.~\ref{fig:mojie_overview} shows the overall interface and main functional modules of the Android application.

\subsection{Offline On-Device Poetry-to-Image Generation}

One of the key features of Mojie is its ability to generate high-quality poetry images entirely in offline and disconnected environments. Once the application and model files are installed on the device, users can perform local inference without requiring network connectivity. This design is particularly valuable for scenarios involving poor connectivity, strict privacy requirements, or a preference for low-latency local execution.

From the user perspective, the interaction is straightforward: a user may input a poem directly or select one from the built-in poetry library, and the application then performs on-device inference to produce a corresponding image. This functionality demonstrates that high-quality poetry-to-image generation can be achieved directly on Android smartphones, validating the feasibility of deploying specialized generative models in practical edge settings. Fig.~\ref{fig:generation_in_mojie} shows an example of offline poetry-to-image generation in Mojie.

\subsection{Built-in Classical Poetry Library}

To improve usability and enrich the application experience, Mojie includes a built-in poetry library containing tens of thousands of classical Chinese poems, covering a wide historical range from the pre-Qin period to the modern era. This extensive collection allows users to conveniently browse poems across dynasties and literary styles within a unified mobile interface.

Beyond simple reading functionality, the poetry library is tightly integrated with the generative model. Users can long-press a poem to directly trigger the image generation process, transforming the selected literary content into a visual result. This interaction design lowers the barrier for non-expert users and provides an intuitive bridge between classical literature and generative AI. Fig.~\ref{fig:built_in_poetry_library} illustrates the built-in poetry library spanning multiple historical periods.

\subsection{Community Sharing and Social Interaction}

In addition to local generation, Mojie incorporates a community-sharing mechanism that supports lightweight social interaction. After generating an image, users can publish the result to the in-app community, where other users can browse the generated content, express appreciation through likes, and participate in discussions through comments.

This design extends the application from a single-user generation tool into a creative and interactive platform. On the one hand, it increases user engagement and encourages repeated use. On the other hand, it creates a channel for sharing diverse visual interpretations of classical poetry, thereby enhancing both the cultural expressiveness and the social value of the system. Fig.~\ref{fig:community_mojie} presents the community page where users can publish and browse poetry-image results.

\subsection{Data Curation}
\label{sec:Data Curation}

The poetry-to-image training corpus is constructed using the dedicated pipeline described in Section~2. In summary, the pipeline transforms approximately 330,000 classical Chinese poems into 1,773,880 high-quality image-text pairs by combining style augmentation, image synthesis, quality filtering, and LLM-based recaptioning.

\subsection{Model Training and Evaluation}
\label{sec:modelTrainingAndEvaluation}

We now present the training and evaluation of the specialized poetry-to-image model deployed in Mojie. The following subsections detail the domain-specific fine-tuning procedure on the high-quality ancient Chinese poetry dataset, the subsequent 4-step distillation strategy for efficient on-device deployment, and both quantitative and qualitative evaluations against mainstream diffusion baselines.

\subsubsection{Training Details}
\label{sec:TrainingDetails}
To adapt our foundational model into a specialized text-to-image generator for ancient Chinese poetry, we fine-tune it on the high-quality ancient Chinese poetry dataset introduced earlier. The entire training process strictly follows the progressive resolution scaling and co-optimization procedures described in the Experiments section. This domain adaptation, inspired by recent advances in personalized and domain-specific diffusion~\cite{ruiz2023dreambooth,gal2022textual,hu2023animate}, aligns the model with the distinctive artistic imagery and visual aesthetics of classical Chinese poetry while preserving its general generative capacity. After foundational fine-tuning, we further apply a 4-step accelerated distillation strategy~\cite{luo2023lcm,luo2023lcm_loRA,hu2024lora} to enable efficient on-device inference. All adaptation experiments are conducted on the Sugon K100 AI cluster using BF16 mixed-precision training.

\subsubsection{Quantitative Results} 
\label{sec:QuantitativeResults} 
To quantitatively evaluate classical Chinese poetry-to-image generation, we assess two complementary aspects: text--image semantic alignment and distributional image fidelity. We compare \textbf{JuZhou~1.0} with representative diffusion baselines, including Stable Diffusion v2.1, SDXL, LCM, SANA-0.6B, and SDXL-Lightning.\footnote{SDXL-Lightning is evaluated with SDXL Base 1.0 as the backbone, and SANA-0.6B uses the 1024-pixel Diffusers checkpoint.}
Since most baselines are primarily optimized for English inputs and provide limited native support for classical Chinese, we evaluate them under both direct Chinese prompting and translated English prompting. In both settings, the original poems are first expanded into generation-oriented prompts using a large language model; for the English setting, the enhanced Chinese prompts are further translated into English before image generation. Chinese-prompt outputs are evaluated using CN-CLIP, while English-prompt outputs are evaluated using CLIP.

\begin{table}[t]
    \centering
    \caption{\textbf{Quantitative evaluation on the poetry image generation benchmark.}
    We report results under Chinese and English prompt settings. Chinese prompts are evaluated with CN-CLIP, while English prompts are evaluated with CLIP. FID is computed against an in-domain held-out reference set.}
    \label{tab:poetry_quantitative}
    \footnotesize
    \setlength{\tabcolsep}{8.0pt}
    \renewcommand{\arraystretch}{1.08}

    \begin{tabular}{lccccccc}
        \toprule
        \textbf{Model} 
        & \textbf{Params$\downarrow$}
        & \textbf{Mobile}
        & \textbf{Prompt}
        & \textbf{Evaluator}
        & \textbf{Steps} 
        & \textbf{CLIP Score$\uparrow$}
        & \textbf{FID$\downarrow$} \\
        \midrule

        \multirow{2}{*}{SDXL}
        & \multirow{2}{*}{2.6B}
        & \multirow{2}{*}{{\color{red}\ding{55}}}
        & CN & CN-CLIP & 50 & 18.74 & 135.65 \\
        & & & EN & CLIP & 50 & 29.96 & \textit{77.04} \\
        \addlinespace[2pt]

        \multirow{2}{*}{SDXL-Lightning}
        & \multirow{2}{*}{2.6B}
        & \multirow{2}{*}{{\color{red}\ding{55}}}
        & CN & CN-CLIP & 4  & 18.54 & 164.62 \\
        & & & EN & CLIP & 4  & 29.21 & 73.70 \\
        \addlinespace[2pt]
    
        \multirow{2}{*}{SDv2.1}
        & \multirow{2}{*}{1.3B}
        & \multirow{2}{*}{{\color{red}\ding{55}}}
        & CN & CN-CLIP & 50 & 17.27 & 109.42 \\
        & & & EN & CLIP & 50 & 28.51 & 77.76 \\
        \addlinespace[2pt]
    
        \multirow{2}{*}{LCM-4-Step}
        & \multirow{2}{*}{\textit{0.86B}}
        & \multirow{2}{*}{{\color{red}\ding{55}}}
        & CN & CN-CLIP & 4  & 19.63 & 159.90 \\
        & & & EN & CLIP & 4  & 28.34 & 98.83 \\
        \addlinespace[2pt]
        
        \multirow{2}{*}{LCM-28-Step}
        & \multirow{2}{*}{\textit{0.86B}}
        & \multirow{2}{*}{{\color{red}\ding{55}}}
        & CN & CN-CLIP & 28 & 18.86 & 159.61 \\
        & & & EN & CLIP & 28 & 27.78 & 93.73 \\
        \addlinespace[2pt]
    
        \multirow{2}{*}{SANA-0.6B}
        & \multirow{2}{*}{\underline{0.6B}}
        & \multirow{2}{*}{{\color{red}\ding{55}}}
        & CN & CN-CLIP & 20 & 22.02 & 84.32 \\
        & & & EN & CLIP & 20 & \underline{29.99} & 75.05 \\
        \midrule
    
        \rowcolor{DeepTeal!15}
        \textbf{JuZhou 1.0 (Ours)}
        & \textbf{0.387B}
        & {\color{green}\ding{51}}
        & CN & CN-CLIP & 28 & \textbf{30.32} & \textbf{26.11} \\
    
        \rowcolor{DeepTeal!15}
        \textbf{JuZhou 1.0 (distilled)}
        & \textbf{0.387B}
        & {\color{green}\ding{51}}
        & CN & CN-CLIP & 4 & 29.01 & \underline{52.54} \\
    
        \bottomrule
    \end{tabular}
    
    \renewcommand{\arraystretch}{1.0}
    \setlength{\tabcolsep}{6pt}
\end{table}

\textbf{Text--Image Alignment.} 
% For the Chinese prompt setting, we use Chinese-CLIP (CN-CLIP) to measure the semantic consistency between generated images and their corresponding enhanced Chinese prompts. As shown in Tab.~\ref{tab:poetry_quantitative}, \textbf{JuZhou~1.0} achieves the highest CN-CLIP score of \textbf{30.32}, outperforming all baselines evaluated with direct Chinese prompts. Its 4-step distilled variant obtains the second-highest score of \underline{29.01}, indicating that DMD2 distillation largely preserves Chinese-language semantic alignment while reducing sampling from 28 to 4 steps. In contrast, the Chinese-prompt scores of general-purpose baselines range from 15.98 to 22.02, suggesting that models primarily optimized for English inputs struggle to interpret classical Chinese poetry without translation.
For the Chinese prompt setting, we use Chinese-CLIP (CN-CLIP) to measure the semantic consistency between generated images and their corresponding enhanced Chinese prompts. As shown in Tab.~\ref{tab:poetry_quantitative}, \textbf{JuZhou~1.0} achieves the highest CN-CLIP score of \textbf{30.32}, substantially outperforming all baselines evaluated with direct Chinese prompts. Its 4-step distilled variant obtains the second-highest score of \underline{29.01}, indicating that DMD2 distillation largely preserves Chinese-language semantic alignment while reducing the sampling process from 28 to 4 steps. In contrast, the Chinese-prompt scores of general-purpose baselines range from 17.27 to 22.02, suggesting that models primarily optimized for English inputs struggle to interpret classical Chinese poetry under direct Chinese prompting.

% For the translated English prompt setting, SDXL achieves the highest CLIP score of \textbf{31.50}, followed by SANA-0.6B with \underline{29.99}, showing that English-oriented models benefit from translation into their preferred input language. However, since CN-CLIP and CLIP use different vision--language encoders and embedding spaces, their scores reflect performance within separate language settings and should not be directly compared on a shared numerical scale.
For the translated English prompt setting, SANA-0.6B achieves the highest CLIP score of \underline{29.99}, closely followed by SDXL with 29.96 and SDXL-Lightning with 29.21. These results show that English-oriented models can benefit from translating Chinese poetry prompts into English. However, since CN-CLIP and CLIP rely on different vision--language encoders and embedding spaces, their scores reflect performance within separate language settings and should not be directly compared on a shared numerical scale.

% \textbf{Distributional Image Fidelity.} 
% We use Fr\'echet Inception Distance (FID)~\cite{fid2023rethinking} to evaluate the distributional similarity between generated images and a held-out reference set from our curated poetry-image benchmark. Since this reference set is derived from the same poetry-oriented data construction pipeline, FID in this benchmark should be interpreted as an in-domain distributional fidelity measure rather than as direct evidence of generalization to arbitrary real-world image distributions. Under this in-domain setting, \textbf{JuZhou 1.0} achieves the lowest FID of \textbf{26.11}, while its 4-step distilled variant obtains the second-lowest FID of \underline{52.54} among all compared settings. Both variants achieve lower FID than the best baseline result, 61.49 from SDXL under the translated English setting. These results indicate that JuZhou 1.0 more closely matches the visual characteristics of the target poetry-image distribution, while the distilled model retains competitive in-domain fidelity despite its substantially reduced sampling cost.

% Because the reference images are derived from the same poetry-oriented data construction pipeline, the reported FID should be interpreted primarily as an in-domain distributional fidelity measure rather than as direct evidence of generalization to arbitrary real-world image distributions. Moreover, FID captures distribution-level similarity and does not independently measure whether each generated image accurately represents the content of its corresponding poem. 

% Overall, the quantitative results demonstrate that JuZhou 1.0 achieves strong native Chinese semantic alignment and in-domain visual fidelity on the constructed classical poetry-to-image benchmark. The distilled variant reduces the sampling process from \textbf{28 steps} to \textbf{4 steps} while retaining a CN-CLIP score of 29.01 and an FID of 52.54, providing a more efficient alternative for deployment-oriented image synthesis.

\textbf{Distributional Image Fidelity.}
We use Fr\'echet Inception Distance (FID)~\cite{fid2023rethinking} to evaluate the distributional similarity between generated images and a held-out reference set from our curated poetry-image benchmark. Because this reference set is constructed using the same poetry-oriented data pipeline, the reported FID should be interpreted primarily as an \emph{in-domain} distributional fidelity measure, rather than as direct evidence of generalization to arbitrary real-world image distributions. Moreover, FID captures distribution-level similarity and does not independently assess whether each generated image accurately reflects the content of its corresponding poem.

Under this in-domain setting, \textbf{JuZhou~1.0} achieves the lowest FID of \textbf{26.11}, while its 4-step distilled variant obtains the second-lowest FID of \underline{52.54} among all compared settings. Both variants substantially outperform the best baseline result, which is 73.70 achieved by SDXL-Lightning under the translated English setting. These results indicate that JuZhou~1.0 more closely matches the visual characteristics of the target poetry-image distribution, while the distilled model retains strong in-domain fidelity despite its substantially reduced sampling cost.

Overall, the quantitative results show that JuZhou~1.0 achieves both strong native Chinese semantic alignment and strong in-domain visual fidelity on the constructed classical poetry-to-image benchmark. Meanwhile, the distilled variant reduces the sampling process from \textbf{28 steps} to \textbf{4 steps} while retaining a CN-CLIP score of 29.01 and an FID of 52.54, providing a substantially more efficient alternative for deployment-oriented image synthesis.

\subsubsection{Qualitative Results}
\label{sec:QualitativeResults}
As illustrated in Fig.~\ref{fig:chinese-poetry}, JuZhou 1.0 demonstrates exceptional qualitative performance in rendering classical Chinese poetry. Unlike conventional baselines that rely on external translation modules---which inherently dilute cultural nuances---JuZhou 1.0 directly and accurately interprets profound literary metaphors and abstract artistic conceptions (e.g., the solitude of a lone boat or the grandeur of an expansive desert). In merely 4 inference steps, the model synthesizes high-fidelity and stylistically diverse artworks, meticulously capturing the intricate aesthetics of traditional Chinese art, from the subtle brushstrokes of ink wash paintings to authentic historical attire and architectural elements. This fusion of poetic abstraction and precise visual grounding highlights the value of native Chinese pretraining for culturally sensitive generative tasks.